\documentclass[11pt]{article}

% Compile from the repository root with pdflatex or tectonic.

\usepackage[margin=1in]{geometry}
\usepackage{amsmath, amssymb}
\usepackage{booktabs}
\usepackage{float}
\usepackage{graphicx}
\usepackage{hyperref}
\usepackage{xcolor}

\graphicspath{{./}{docs/figures/}{figures/}}

\hypersetup{
  colorlinks=true,
  linkcolor=blue,
  urlcolor=blue,
  citecolor=blue
}

\setlength{\parindent}{0pt}
\setlength{\parskip}{0.65em}
\setlength{\emergencystretch}{2em}

\newcommand{\reportfigure}[3]{%
  \begin{figure}[H]
    \centering
    \includegraphics[width=0.98\linewidth]{#1}
    \caption{#2}
    \label{#3}
  \end{figure}
}
\newcommand{\reportprob}{q}

\title{\textbf{Observer--Actor Dynamic: Technical Report}}
\author{}
\date{}

\begin{document}

\maketitle

\begin{abstract}
This report presents controlled experiments for the Observer--Actor learning
framework. Each experiment separates one learned part of the system and records
both aggregate reward and the specific quantity that part is supposed to learn.

The first two experiments study reward-weighted policy learning in a single fixed
state and in multiple fully observed states. The one-to-one imperfect-state
experiment starts the imperfect-state track:
an Actor-compressor observes the state and sends a scalar message, while an
Observer acts from that message. The imperfect-reward experiments study a
single Observer with a limited budget for following many Actors, followed by
the group-level version with many Observers and imperfect reward functions. The
imperfect-state experiments then scale the state-information problem: one
Observer selects among several Actor-compressors, and then several Observers
each select Actor-compressors and act from their messages.
\end{abstract}

\tableofcontents

\section{Overview}

The paper studies a setting where the acting object and the information object
can be different. We call the acting object the Actor. We call the information
object the Observer. The Actor chooses an action from a policy
\(\pi_\theta\). The Observer can affect learning either by reporting reward
information or by sending state information.

There are two main information problems in the report. In the imperfect-reward
cases, each Actor can see the state and choose actions, but reward feedback
comes through Observer following. The learned quantities are the following
rates \(\mu\), estimated Actor utilities \(\hat U\), and in the joint cases
the Actor policy. In the imperfect-state cases, the direction is different.
An Actor-compressor sees state information and sends a message
\(m=f_\phi(s)\). The Observer receives the message, chooses an environmental
action, and sends reward feedback. In the multi-compressor cases the Observer
also learns which Actor-compressor to listen to. Those compressor-selection
probabilities are denoted by \(\lambda\).

The experiments are built as a ladder. First we test plain policy learning.
Then we add scalar state compression. Then we add selection among
Actor-compressors. Finally, we add several Observers, each with its own
selection rule and policy. This order is intended so that, when a joint
experiment learns poorly, we can tell whether the problem is policy learning, message
compression, compressor selection, or the coupling between them.

For that reason the report tracks more than reward. Reward says whether the
system performs well; the learned mappings show which part moved. We also track
message values, compression weights, following rates,
compressor-selection probabilities, and per-state action probabilities. A
result is treated as reliable when both reward and the relevant learned
quantity move in the expected direction.

\section{Method}

This section fixes the common notation used in the experiments. An Actor policy
maps its input to action logits. The softmax of those logits gives the action
probabilities. In a few controlled experiments, the Actor policy is fixed so that
only the Observer, following rate, or selection rule is being tested. When the
policy is learned, it is updated from sampled rewards unless the text explicitly
marks the run as a controlled comparison.

For the action-learning cases, the training loss is
\[
  \mathcal{L}_t
  =
  -\hat r_t \log \pi_\theta(a_t \mid x_t),
\]
where \(\hat r_t\) is the Observer-reported reward, \(x_t\) is the Actor input,
and \(a_t\) is the sampled action. Minimizing this loss performs stochastic
gradient ascent on expected reward:
\[
  \theta_{t+1}
  \leftarrow
  \theta_t
  +
  \alpha r_t \nabla_\theta \log \pi_\theta(a_t \mid x_t).
\]

The policy input differs by case:
\[
  x_t =
  \begin{cases}
    s_{\mathrm{static}}, & \text{single-state action learning},\\
    s_t, & \text{state-dependent action learning},\\
    m_t=f_\phi(s_t), & \text{one-to-one scalar-compression experiments},\\
    m_t=f_{\phi_{k_t}}(s_t), & \text{multi-compressor imperfect-state experiments}.
  \end{cases}
\]

For the imperfect-state experiments, the same sampled reward can update three
objects. If compressor \(k_t\) is selected and sends
\(m_t=f_{\phi_{k_t}}(s_t)\), the policy update is
\[
  \theta_{t+1}
  \leftarrow
  \theta_t
  +\alpha_\theta r_t
  \nabla_\theta\log\pi_\theta(a_t\mid m_t).
\]
When compression is learned, the selected compressor receives the gradient
through the policy:
\[
  \phi_{k_t,t+1}
  \leftarrow
  \phi_{k_t,t}
  +\alpha_\phi r_t
  \nabla_m\log\pi_\theta(a_t\mid m_t)
  \nabla_{\phi_{k_t}}f_{\phi_{k_t}}(s_t).
\]
When selection is learned, its logits receive
\[
  w_{t+1}
  \leftarrow
  w_t+\alpha_w r_t\nabla_w\log\lambda_{k_t}(w_t).
\]
Each subsection states which of \(\phi\), \(w\), and \(\theta\) is learned and
which is fixed.

The imperfect-state experiments are divided into stages. First, in the one-to-one
case, one Actor-compressor learns a scalar message for one Observer with a
fixed policy. Next, one Observer selects among several Actor-compressors while
the selected compressor learns its message. Finally, several Observers each
select Actor-compressors and act from the received messages. In the learned
policy experiments, the Observer policy parameters \(\theta\), the
compressor-selection logits \(w\), and the selected compression parameters
\(\phi\) can be fixed or learned depending on which part is being tested.

\textbf{Distance-based softmax policy for scalar messages.}
The learned-policy imperfect-state experiments use a scalar message \(m\). A plain
linear softmax policy has trouble when one scalar line must be split into many
action regions. The policy used in the current experiments keeps the scalar message,
but scores each action by distance from the message to that action's center:
\[
  \ell_a(m)=-\beta(m-a)^2,\qquad
  \pi_\theta(a\mid m)=\operatorname{softmax}(\ell(m))_a .
\]
Here \(\ell_a(m)\) is the logit for action \(a\), and \(\beta\) controls how
sharp the action preference is. For example, if \(K=8\), \(m=5.1\), and
\(\beta=32\), action 5 receives the largest logit because 5 is the closest
action center. If state 5 emits \(m=0.2\), the score for the correct action has
derivative \(-2\beta(m-5)\), which pushes the message back toward 5 during joint
compression and policy learning.

When there are \(K\) states and \(A<K\) actions, neighboring states share an
action. We first map the message onto the action scale,
\[
  z={A\over K}m+\left({A\over 2K}-{1\over 2}\right),
  \qquad
  \ell_a=-\beta(z-a)^2.
\]
For example, with \(K=128\) states and \(A=32\) actions, each action covers four
neighboring states.

The learned-policy scaling experiments use 300k steps,
\(\epsilon_w=\epsilon_a=0.05\), \(\alpha_\phi=10^{-6}\),
\(\alpha_w=10^{-4}\), \(\alpha_\theta=3\cdot10^{-5}\), and
\(\beta=32\), except where a table states another value. Here
\(\epsilon_w\) is the uniform compressor-selection floor and
\(\epsilon_a\) is the uniform action-sampling floor:
\[
  \widetilde{\pi}(a\mid m)
  =(1-\epsilon_a)\pi_\theta(a\mid m)+{\epsilon_a\over A}.
\]
The learned distributions, rather than the uniform mixtures, provide the
log-probabilities in the reward-gradient updates.

The imperfect-reward experiments use a separate setup: the state is fully
observed, and the learned objects are budgeted following rates \(\mu\), utility
estimates \(\hat U\), and in the joint case Actor policies.

For the first two action-learning experiments, the default learning-rate grid is
\[
  \{0.001,\;0.01,\;0.05,\;0.1\},
\]
with 5 random seeds and 5000 training steps. Later experiments state their own
learning rates, random seeds, and training steps. The main reward metric is
expected reward under the current policy. The figures also report the learned
quantity required by the case:
action probabilities in the action-learning experiments, message behavior in the
one-to-one imperfect-state experiment,
compressor-selection probabilities and correct-action probability in the
imperfect-state experiments, and following rates in the imperfect-reward
experiments.

The tables use consistent metric definitions. Learning rate is the optimizer
step size. Expected reward is the reward the current stochastic policy would
receive on average before sampling an action. In one state this is
\(\sum_a\pi(a\mid x)r(a)\); across states it is also averaged over the state
distribution. End expected reward is this quantity at the last reported
training step, averaged over seeds. End sampled reward averages the last 100
rewards actually sampled during training. High-reward action probability is
the total probability assigned to the maximum-reward action after averaging
over the relevant states and selected Observers. Success rate is the fraction
of seeds reaching at least 95\% of the reward ceiling. Plots show trajectories;
tables summarize end-of-run behavior.

For example, suppose one state has rewards \(r(0)=0\), \(r(1)=1\), and
\(r(2)=0.2\), and the current policy assigns probabilities
\((0.05,0.90,0.05)\). The expected reward is
\(0.05\cdot0+0.90\cdot1+0.05\cdot0.2=0.91\), even though the next sampled
reward will be one of \(0\), \(1\), or \(0.2\). In a three-seed run, an
end expected reward of \(0.96\) means the three end expected rewards average to
\(0.96\). If the reward ceiling is \(1.0\), a seed is counted as successful
once it reaches at least \(0.95\). If high-reward action probability is
\(0.90\), the current policy assigns 90\% probability to the best action after
receiving the state input or selected compressor message.

\section{Single-State Action Learning}

This case tests the basic policy-gradient update. There is one fixed state
\(s_{\mathrm{static}}\), three actions, and the reward table is
\[
  r(0)=0.0,\qquad r(1)=1.0,\qquad r(2)=0.2.
\]
The only learned object is the Actor's action distribution in this one state.
The objective is
\[
  J(\theta)
  =
  \mathbb{E}_{a \sim \pi_\theta(\cdot \mid s_{\mathrm{static}})}[r(a)]
  =
  \sum_a \pi_\theta(a \mid s_{\mathrm{static}}) r(a).
\]
The optimal behavior is to put almost all probability on action 1.

\reportfigure{report_case1_reward_curves}
  {Single-state reward learning curves. Expected reward gives the policy-level
  performance, while average sampled reward shows the recent reward actually
  experienced during training. Curves show the mean over 5 seeds for each
  learning rate.}
  {fig:case1-reward-curves}

\reportfigure{report_case1_action_probabilities}
  {Single-state action probabilities. The probability of the reward-maximizing action rises
  toward one, while the lower-reward actions lose probability. Curves show the
  mean over 5 seeds at learning rate \(0.05\).}
  {fig:case1-action-probs}

\begin{table}[H]
\centering
\small
\begin{tabular}{cccc}
\toprule
Learning rate & End expected & End sampled & Success rate \\
\midrule
0.001 & 0.772 & 0.795 & 0.00 \\
0.01  & 0.992 & 0.989 & 1.00 \\
0.05  & 0.999 & 0.996 & 1.00 \\
0.1   & 0.999 & 0.998 & 1.00 \\
\bottomrule
\end{tabular}
\caption{Single-state action-learning end-of-run summary over 5 seeds.}
\label{tab:case1-summary}
\end{table}

This shows the intended update directly: the learned object is the action
distribution, and the probability of the best action increases.

\section{State-Dependent Action Learning}

The state-dependent action-learning experiment tests whether the same update can
learn different actions in different
fully observed states. There are three states, sampled uniformly, and three
actions. States are represented as state vectors. The reward structure is:
\[
  \begin{array}{c|ccc}
  & a=0 & a=1 & a=2 \\
  \hline
  s=0 & 1.0 & 0.0 & 0.2 \\
  s=1 & 0.1 & 1.0 & 0.0 \\
  s=2 & 0.0 & 0.2 & 1.0
\end{array}
\]

The learned object is now a conditional policy: for each state, the Actor
should put high probability on that state's best action.

The objective is
\[
\begin{aligned}
  J(\theta)
  &=
  \mathbb{E}_{s \sim D}
  \left[
    \mathbb{E}_{a \sim \pi_\theta(\cdot \mid s)}[r(s,a)]
  \right]
  \\
  &=
  \sum_s D(s) \sum_a \pi_\theta(a \mid s) r(s,a).
\end{aligned}
\]

\reportfigure{report_case2_reward_curves}
  {State-dependent reward learning curves. Expected reward gives the policy-level
  performance across the uniformly sampled states, while average sampled reward
  shows the recent reward actually experienced during training. Curves show the
  mean over 5 seeds for each learning rate.}
  {fig:case2-reward-curves}

\reportfigure{report_case2_state_probabilities}
  {State-conditioned action probabilities. Each curve shows the probability assigned to the correct
  action for one state. The model learns the three state-action relationships
  while average reward improves. Curves show the mean over 5 seeds
  at learning rate \(0.05\).}
  {fig:case2-state-probs}

\begin{table}[H]
\centering
\small
\begin{tabular}{cccc}
\toprule
Learning rate & End expected & End sampled & Success rate \\
\midrule
0.001 & 0.499 & 0.504 & 0.00 \\
0.01  & 0.939 & 0.939 & 0.00 \\
0.05  & 0.991 & 0.986 & 1.00 \\
0.1   & 0.996 & 0.991 & 1.00 \\
\bottomrule
\end{tabular}
\caption{State-dependent action-learning end-of-run summary over 5 seeds.}
\label{tab:case2-summary}
\end{table}

The per-state probability plot shows the internal behavior that the scalar reward
curve hides: each conditional action distribution learns its own best action.

\section{Imperfect State with One Actor-Compressor and One Observer}

The one-to-one imperfect-state experiment is the first imperfect-state experiment.
It follows the updated paper
direction for the state-information side: an Actor-compressor observes the
state and sends a message, and an Observer acts from that message. This case
has one Actor-compressor and one Observer. The three experiments learn
compression \(\phi\) alone, policy \(\theta\) alone, and then
\(\phi\) and \(\theta\) together.

The controlled state space has \(K\) states, sampled uniformly. State \(k\) is
represented by a cumulative binary vector with \(k\) ones:
\[
\begin{array}{rcl}
  s_0 &=& [0,0,\ldots,0],\\
  s_1 &=& [1,0,\ldots,0],\\
  s_2 &=& [1,1,0,\ldots,0],\\
      &\vdots& \\
  s_{K-1} &=& [1,1,\ldots,1,0].
\end{array}
\]
The intended scalar message is the sum of the vector:
\[
  m=f_\phi(s_k)\approx \sum_{\ell=1}^K s_{k,\ell}=k.
\]
For example, with \(K=5\), state
\([1,1,1,0,0]\) has sum \(3\), so the useful message is near \(m=3\).
The correct action is the state index,
\[
  a^\star(s_k)=k.
\]
Correct actions receive reward \(1.0\), and wrong actions receive \(-0.1\).

\subsection{Compression Only: Learn \(\phi\), Fixed Observer Policy}

The compressor is linear:
\[
  f_\phi(s)=u^\top s+b,
\]
where \(u\) and \(b\) are learned. The Observer policy is fixed in this
compression-only experiment. It prefers action \(a\) when the scalar message is near
that action index:
\[
  \ell_a(m)=-\beta(m-a)^2,\qquad
  \pi(a\mid m)=\operatorname{softmax}(\ell(m))_a .
\]
This fixed policy gives \(\phi\) a clear target: state \(k\) should produce a
message near \(k\).

The main internal quantities are message mean absolute error and message
\(R^2\). The mean absolute error compares \(f_\phi(s_k)\) with \(k\). The
\(R^2\) value measures whether the learned messages line up with the true state
IDs up to an affine rescaling. A value near \(1.0\) means the messages have the
right ordering and nearly linear spacing. ``Missed states'' counts states whose
correct-action probability is below \(0.95\).

\reportfigure{updated_section_5_1_scalar_sum_k64_long_curves}
  {One-to-one scalar-compression learning at \(K=64\). The reward and
  correct-action probability rise when the compression learning rate is in the
  stable range. Very small learning rates move slowly, while overly large rates
  can move the scalar messages too aggressively.}
  {fig:case3-scalar-k64}

\reportfigure{updated_section_5_1_scalar_sum_k64_long_final_weights}
  {One-to-one end-of-run compression plot at \(K=64\). The learned scalar messages are
  plotted against the target state IDs. A near-straight increasing line means
  the compressor has learned the intended sum-like state message.}
  {fig:case3-scalar-k64-weights}

\begin{table}[H]
\centering
\small
\resizebox{\linewidth}{!}{%
\begin{tabular}{cccccccc}
\toprule
\(K\) & Steps & Best lr & End expected & Min correct & Missed states & Message MAE & Message \(R^2\) \\
\midrule
3   & 30k  & 0.10  & 0.987 & 0.965 & 0.00  & 0.555 & 1.000 \\
5   & 30k  & 0.05  & 0.977 & 0.965 & 0.00  & 0.277 & 0.988 \\
8   & 30k  & 0.05  & 0.971 & 0.965 & 0.00  & 0.183 & 0.991 \\
16  & 30k  & 0.02  & 0.966 & 0.965 & 0.00  & 0.052 & 0.999 \\
32  & 120k & 0.005 & 0.963 & 0.965 & 0.00  & 0.020 & 1.000 \\
64  & 240k & 0.005 & 0.962 & 0.965 & 0.00  & 0.010 & 1.000 \\
\bottomrule
\end{tabular}
}
\caption{One-to-one scalar-compression scaling over 5 seeds. The fixed Observer
policy has a finite sharpness, so the practical reward ceiling is below
\(1.0\).}
\label{tab:case3-scalar-scaling}
\end{table}

Compression-only learning works through \(K=64\). The learned scalar messages
line up with the state IDs and the fixed policy selects the correct action with
high probability.

\subsection{Policy Only: Learn \(\theta\), Fixed Correct Message}

This experiment learns the Observer policy \(\theta\) while the message is fixed
to the correct scalar value \(m=k\). The Observer receives the right state
message and only has to learn which action to take from that message. The policy
uses the distance-based softmax form defined in the Method section.

\begin{table}[H]
\centering
\small
\begin{tabular}{ccccc}
\toprule
\(K,A\) & Reward & Correct-action prob. & Msg MAE & Missed states \\
\midrule
128, 32 & 1.000 & 1.000 & 0.000 & 0.00 \\
\bottomrule
\end{tabular}
\caption{One-to-one policy-only experiment. The message is fixed to the correct
state value and only the Observer policy is learned.}
\label{tab:one-to-one-policy-only}
\end{table}

\subsection{Compression and Policy: Learn \(\phi,\theta\), Single Message Path}

This experiment learns both the scalar message \(f_\phi(s)\) and the Observer
policy \(\pi_\theta(a\mid m)\). There is only one Actor-compressor and one
Observer, so the message path is fixed and the selection variable is absent.

\begin{table}[H]
\centering
\small
\begin{tabular}{ccccc}
\toprule
\(K,A\) & Reward & Correct-action prob. & Msg MAE & Missed states \\
\midrule
128, 32 & 0.999 & 0.971 & 0.146 & 0.00 \\
\bottomrule
\end{tabular}
\caption{One-to-one compression-and-policy experiment. The scalar message and
the Observer policy are learned together on the single message path.}
\label{tab:one-to-one-phi-theta}
\end{table}

With the distance-based softmax policy, the scalar learned-policy case works at
\(K=128,A=32\), including the case where \(\phi\) and \(\theta\) learn
together. This keeps the scalar message setup and uses a policy shape that
matches the intended sum message.

\section{Imperfect Reward: One Observer and Many Actors}

This section corresponds to the paper's one-Observer imperfect-reward setting.
Every Actor observes the full state. After an Actor acts, the Observer sees the
reward for that action and reports it with probability determined by the
following rate. A single Observer watches \(N\) Actors and has a limited budget
for deciding which Actors to follow.

A useful way to read this section is as one instructor watching several
students. The students are the Actors: they produce answers. The instructor is
the Observer: it estimates which students tend to produce higher-reward
answers and spends more of its limited following budget on those students.
This picture is only an intuition, but it matches the roles in the update:
\(\hat U_j\) estimates Actor \(j\)'s average quality, and \(\mu_j\)
controls how often the Observer gives that Actor reward feedback.

The new learnable object is the budgeted following-rate vector
\[
  \mu = (\mu_1,\dots,\mu_N),\qquad \mu_j \ge 0,\qquad \sum_{j=1}^{N}\mu_j \le B,
\]
where \(\mu_j\) is how strongly the Observer follows Actor \(j\). The following
rate sets the probability that the Observer actually reports the reward it saw,
\[
  \reportprob(\mu_j) = 1 - e^{-c\,\mu_j},\qquad c>0,
\]
so larger attention means more frequent feedback. The Observer also keeps a
running utility estimate \(\hat U_j\) of each Actor's average reward. The
tracked variables for this case are \(\mu_j(t)\) and \(\hat U_j(t)\): the main
experiment is reward improvement together with Observer reallocation toward the
genuinely higher-utility Actors.

\textbf{Setup conventions.} The controlled tests use a single state (so each
fixed Actor's utility is known analytically), three actions, scalar state-ID
inputs mapped through a learnable embedding, and a hard-coded reward table.
One Actor \(j_t\sim p\) acts per step with \(p\) uniform; the Observer sees
only the realized reward of the acting Actor. In the fixed-policy experiment,
the Actor policies are fixed and only the Observer's \(\mu\) and \(\hat U\)
are learned. In the joint-learning experiment, each Actor additionally learns
its softmax policy.

\textbf{Updates.} When Actor \(j\) acts, the utility estimate uses a fast
stochastic-approximation step
\[
  \hat U_{j,t+1} = \hat U_{j,t} + \alpha_U\bigl(r_t - \hat U_{j,t}\bigr),
\]
while the following rates take a slower projected-gradient-ascent step on
\(\sum_j \reportprob(\mu_j)\hat U_j\),
\[
  \hat\mu_{j}(t+1) = \mu_j(t) + \alpha_\mu\,\reportprob'(\mu_j(t))\,\hat U_{j,t},
  \qquad
  \mu(t+1) = \Pi_{\{\mu\ge 0,\ \sum_j\mu_j\le B\}}\bigl(\hat\mu(t+1)\bigr),
\]
with \(\reportprob'(\mu)=c\,e^{-c\mu}\). The projection has a closed form via the
capped-simplex projection. Two-timescale behavior is obtained by choosing
\(\alpha_U \gg \alpha_\mu\), so the utilities are accurate before the attention
moves.

\subsection{Learnable Following Rates with Fixed Actor Policies}

This subsection studies the following-rate update alone. With fixed Actors of
known utility, projected gradient ascent should move the budget \(B\) toward
the best Actors.

\textbf{Controlled test 1: one strictly-best Actor.} A single state, three fixed
Actors, and reward row \(R=[0.9,0.5,0.1]\). Each fixed Actor uses a sharp
softmax centered on its target action, so the analytical utilities round to
\([0.9,0.5,0.1]\). We
use \(B=1\), \(c=2\), \(\alpha_U=0.10\), \(\alpha_\mu=0.02\), and 4000 steps. The
expected behavior is \(\hat U\to[0.9,0.5,0.1]\), \(\mu\) concentrating on Actor
0, the low-utility Actor 2 falling to \(\mu_2\approx 0\), and the budget
binding.

\reportfigure{paper4_1_single_best}
  {Single-best following-rate result. Left: following rates \(\mu_j(t)\) ---
  the high-utility Actor 0 gains budget while the low-utility Actor 2 drops to
  zero. Right: utility estimates \(\hat U_j(t)\) match the dotted analytical
  utilities. Mean over 5 seeds.}
  {fig:section4-single-best}

\begin{table}[H]
\centering
\small
\begin{tabular}{lccc}
\toprule
Quantity & Actor 0 & Actor 1 & Actor 2 \\
\midrule
Analytical utility \(U_j\)        & 0.900 & 0.500 & 0.100 \\
Final estimate \(\hat U_j\)       & 0.900 & 0.500 & 0.100 \\
Final following rate \(\mu_j\)    & 0.647 & 0.353 & 0.000 \\
\bottomrule
\end{tabular}
\caption{Single-best following-rate test, mean over 5 seeds. The budget binds
(\(\sum_j\mu_j=1.000\)) and attention is ordered by true utility
(\(\mu_0>\mu_1>\mu_2\)).}
\label{tab:section4-single-best}
\end{table}

\textbf{Controlled test 2: tie.} The same setup with \(R=[0.9,0.9,0.1]\), so
Actors 0 and 1 are equally useful and Actor 2 has low utility. The expected
behavior is that the budget splits across the tied optimal Actors while Actor 2
loses attention.

\reportfigure{paper4_1_tie}
  {Following-rate tie result. The two equally useful Actors share the budget
  symmetrically and the low-utility Actor 2 drops to zero. Because the two top
  Actors are genuinely symmetric, the across-seed mean equals the per-seed
  behavior. Mean over 5 seeds.}
  {fig:section4-tie}

\begin{table}[H]
\centering
\small
\begin{tabular}{lccc}
\toprule
Quantity & Actor 0 & Actor 1 & Actor 2 \\
\midrule
Analytical utility \(U_j\)        & 0.900 & 0.900 & 0.100 \\
Final estimate \(\hat U_j\)       & 0.900 & 0.900 & 0.100 \\
Final following rate \(\mu_j\)    & 0.500 & 0.500 & 0.000 \\
\bottomrule
\end{tabular}
\caption{Following-rate tie test, mean over 5 seeds. Mass leaves the low-utility
Actor and splits evenly across the tied optimal set.}
\label{tab:section4-tie}
\end{table}

Both tests show the learned variables directly: the utility estimates recover
the true utilities and the following rates move to the high-utility set under
the budget constraint.

\subsection{Learnable Following Rates and Actor Policies}

Here each Actor also learns. A reported step --- one where the Bernoulli draw
\(\mathrm{Bernoulli}(\reportprob(\mu_j))\) succeeds --- drives a REINFORCE update
\[
  \theta_j \leftarrow \theta_j
  + \alpha_\theta\, r_t\, \nabla_{\theta_j}\log\pi_{\theta_j}(a_t\mid s_t).
\]
The Observer's \(\mu\) and \(\hat U\) update as in the fixed-policy experiment.

\textbf{Setup.} Three states (uniform), three actions, three Actors, reward
table
\[
  R=\begin{pmatrix} 1.0 & 0.1 & 0.2 \\ 0.1 & 1.0 & 0.0 \\ 0.0 & 0.2 & 1.0 \end{pmatrix},
\]
optimal actions \([0,1,2]\), \(B=3\), \(c=2\) (so the initial \(\mu=[1,1,1]\)
gives \(\reportprob(1)\approx 0.865\), meaning most acted steps are reported),
\(\alpha_U=0.20\),
\(\alpha_\mu=0.01\), \(\alpha_\theta=0.20\), and 30000 steps. Since every Actor
can become optimal and the budget is large enough to fund all of them, the
correct fixed point is each Actor's optimal-action probability approaching one,
\(\hat U_j\to 1\), and a symmetric allocation \(\mu\to[1,1,1]\).

\reportfigure{paper4_2_joint}
  {Joint following-rate and Actor-policy learning. Left: per-Actor optimal-action
  probability climbs to \(\approx 0.999\). Middle: following rates \(\mu_j(t)\)
  remain near the feasible symmetric allocation. Right: utility estimates
  \(\hat U_j(t)\) rise toward the per-Actor ceiling of one. Mean over 5 seeds.}
  {fig:section4-joint}

\begin{table}[H]
\centering
\small
\begin{tabular}{lccc}
\toprule
Quantity & Actor 0 & Actor 1 & Actor 2 \\
\midrule
Final optimal-action prob.        & 0.999 & 0.999 & 0.999 \\
Final following rate \(\mu_j\)    & 0.999 & 1.001 & 1.000 \\
\bottomrule
\end{tabular}
\caption{Joint following-rate and Actor-policy learning, mean over 5 seeds.
All Actors reach
near-deterministic optimal policies and the budget stays feasible
(\(\sum_j\mu_j=3.000\)).}
\label{tab:section4-joint}
\end{table}

In this controlled setup, joint learning is stable with the two-timescale choice
(\(\alpha_U\gg\alpha_\mu\)): the fast utility estimate tracks the improving
Actors, and the slow following-rate update keeps \(\mu\) feasible throughout.
The figures show the learned quantities directly: policies improve, utilities
rise, attention stays within budget, and end-of-run reward is high.

\section{Group-Level Interactions: Imperfect Reward}

This section is the paper's group-level imperfect-reward extension of the
previous one. Instead of a single Observer watching \(N\) Actors, there are now
\(N\) Observers
\(O_1,\dots,O_N\) and \(N\) Actors \(A_1,\dots,A_N\). Every Actor still sees the
state perfectly, but each Observer \(O_i\) now has its \emph{own} reward function
\(R_i\): when the active Actor performs an action, Observer \(i\) perceives
reward \(R_i(s_t,a_t)\). Because the \(R_i\) differ, two Observers may value the
same Actor differently, and each Observer maintains its own budgeted
following-rate vector \(\mu_i=(\mu_{i,1},\dots,\mu_{i,N})\) and utility estimate
\(\hat U_{i,j}\).

\textbf{Setup conventions.} As before, the controlled fixed-policy test uses a
single state, three actions, scalar state-ID inputs mapped through a learnable
embedding, and hard-coded per-Observer reward tables; one Actor \(j_t\sim p\)
acts per step with \(p\) uniform. Each Observer maintains its own
capped-simplex following rates with budget \(B\) and reporting probability
\(\reportprob(\mu_{i,j})=1-e^{-c\,\mu_{i,j}}\).

\textbf{Updates.} When Actor \(j_t\) acts, every Observer updates its utility
estimate for that Actor on the fast timescale using its own perceived reward,
\[
  \hat U_{i,j_t,t+1} = \hat U_{i,j_t,t}
  + \alpha_U\bigl(R_i(s_t,a_t) - \hat U_{i,j_t,t}\bigr),
\]
and updates its following rates by the same projected gradient ascent as the
single-Observer case. In the joint-learning setting, Actor \(j\) additionally
learns from an Observer-aggregated reward (paper Eq.~29),
\[
  r_{\mathrm{eff}}(j) = \sum_{i=1}^{N}\reportprob(\mu_{i,j})\,R_i(s_t,a_t),
  \qquad
  \theta_j \leftarrow \theta_j
  + \alpha_\theta\, r_{\mathrm{eff}}(j)\,
  \nabla_{\theta_j}\log\pi_{\theta_j}(a_t\mid s_t),
\]
so an Actor's learning signal reflects both the reward for its action and how
often Observers report that reward.

\subsection{Multiple Observers and Actors with Learnable Following Rates}

This subsection studies whether each Observer, using its own imperfect reward
function, reallocates its budget toward the Actors \emph{it} values most while
Actor policies are fixed.

\textbf{Controlled test: assortative reward.} A single state, three fixed Actors
(Actor \(j\) pinned to action \(j\)), and three Observers with diagonal reward
tables \(R_0=[0.9,0.4,0.1]\), \(R_1=[0.1,0.9,0.4]\), \(R_2=[0.4,0.1,0.9]\), so
Observer \(i\) most values Actor \(i\). We use \(B=2\), \(c=1\),
\(\alpha_U=0.10\), \(\alpha_\mu=0.05\), and 8000 steps. With these utilities the
constrained optimum is the closed-form KKT allocation \(\mu=[1.405,0.595,0]\)
(reordered per Observer): the highest-utility Actor receives the most budget and
the worst Actor drops to zero.

\reportfigure{paper6_1_assortative}
  {Assortative following with heterogeneous Observers. Each panel shows one
  Observer's following rates \(\mu_{i,j}(t)\); every Observer concentrates its
  budget on the Actor it values most, so the three Observers end with different
  Actor preferences according to their own reward functions.
  Mean over 5 seeds.}
  {fig:section6-assortative}

\begin{table}[H]
\centering
\small
\begin{tabular}{lccc}
\toprule
Final following rate \(\mu_{i,j}\) & Actor 0 & Actor 1 & Actor 2 \\
\midrule
Observer 0 (\(U=[0.9,0.4,0.1]\)) & 1.405 & 0.595 & 0.000 \\
Observer 1 (\(U=[0.1,0.9,0.4]\)) & 0.000 & 1.405 & 0.595 \\
Observer 2 (\(U=[0.4,0.1,0.9]\)) & 0.595 & 0.000 & 1.405 \\
\bottomrule
\end{tabular}
\caption{Assortative following test, mean over 5 seeds. Each Observer's budget
binds (\(\sum_j\mu_{i,j}=2.000\)) and concentrates on its highest-utility Actor,
matching the analytical KKT allocation.}
\label{tab:section6-assortative}
\end{table}

The empirical following rates match the analytical allocation, and the three
Observers move toward three different Actors. This is the group-level experiment:
different reward functions produce different budget allocations.

\subsection{Multiple Observers and Actors with Joint Learning}

Here each Actor also learns, using the Observer-aggregated reward
\(r_{\mathrm{eff}}(j)=\sum_i\reportprob(\mu_{i,j})R_i\) defined above.

\textbf{Setup.} Three states (uniform), three actions, three Actors, and three
Observers sharing the diagonal reward table
\[
  R=\begin{pmatrix} 1.0 & 0.1 & 0.2 \\ 0.1 & 1.0 & 0.0 \\ 0.0 & 0.2 & 1.0 \end{pmatrix},
\]
optimal actions \([0,1,2]\), \(B=3\), \(c=2\), \(\alpha_U=0.20\),
\(\alpha_\mu=0.01\), \(\alpha_\theta=0.10\), and 30000 steps. Because all Actors
can become optimal, they become equally valuable to every Observer, so the
correct fixed point is each Actor's optimal-action probability approaching one
and a symmetric allocation \(\mu_i\to[1,1,1]\) for every Observer.

\reportfigure{paper6_2_joint}
  {Group-level joint learning. Left: per-Actor optimal-action probability climbs
  to \(\approx 0.999\) under the Observer-aggregated reward. Middle: Observer~0's
  following rates \(\mu_{0,j}(t)\) stay near the feasible symmetric allocation.
  Right: Observer~0's utility estimates \(\hat U_{0,j}(t)\) rise toward the
  per-Actor ceiling of one. Mean over 5 seeds.}
  {fig:section6-joint}

\begin{table}[H]
\centering
\small
\begin{tabular}{lccc}
\toprule
Quantity & Actor 0 & Actor 1 & Actor 2 \\
\midrule
Final optimal-action prob.            & 0.999 & 0.999 & 0.999 \\
Final following rate \(\mu_{i,j}\) (any \(i\)) & 1.000 & 1.000 & 1.000 \\
\bottomrule
\end{tabular}
\caption{Group-level joint learning, mean over 5 seeds. All Actors reach
near-deterministic optimal policies under the aggregated reward, and every
Observer's budget stays feasible (\(\sum_j\mu_{i,j}=3.000\)).}
\label{tab:section6-joint}
\end{table}

In this controlled setup, joint group learning is stable with the two-timescale
choice. The aggregated reward \(\sum_i\reportprob(\mu_{i,j})R_i\) trains the
Actors, and every Observer's following rates stay within budget.
\section{Imperfect State with Multiple Actor-Compressors and One Observer}

This section is the one-Observer, many-compressor imperfect-state setting. One
Observer chooses which Actor-compressor to listen to, receives that scalar
message, and then acts. The three Actor-compressors are:
\[
\begin{array}{c|c|c}
\text{Compressor} & \text{state input} & \text{role} \\
\hline
0 & \text{zero vector} & \text{uninformative baseline} \\
1 & \text{first coordinate only} & \text{partial state information} \\
2 & \text{full cumulative state vector} & \text{can identify the state}.
\end{array}
\]
The Observer has selection logits \(w\), with selection probabilities
\[
  \lambda_i(w)=\frac{\exp(w_i)}{\sum_\ell \exp(w_\ell)}.
\]
During training we sample from a distribution with a small uniform floor,
\[
  \tilde\lambda_i=(1-\epsilon_w)\lambda_i+\frac{\epsilon_w}{N},
\]
where \(N\) is the number of Actor-compressors. This keeps every compressor
sampled sometimes, while the reported \(P(\mathrm{full})\) remains the learned
probability \(\lambda_2\).

\subsection{Compression Only: Learn \(\phi\), Fixed Full Selection and Fixed Policy}

This experiment fixes the Observer to listen to the full-information compressor
and fixes the Observer policy. The only learned object is the full compressor's
message rule \(f_{\phi_2}\).

\begin{table}[H]
\centering
\small
\begin{tabular}{ccccc}
\toprule
\(K\) & Reward & \(P(\mathrm{full})\) & Correct-action prob. & Missed@0.90 \\
\midrule
64 & 0.962 & 1.000 & 0.965 & 0.00 \\
\bottomrule
\end{tabular}
\caption{One-to-many compression-only experiment. Selection and policy are
fixed, and only the full compressor message is learned.}
\label{tab:one-to-many-phi-only}
\end{table}

\subsection{Selection Only: Learn \(w\), Fixed Correct Messages and Fixed Policy}

This experiment fixes all compressor messages and the Observer policy. The full
compressor sends the correct scalar state message, while the other compressors
carry incomplete information. The only learned object is the selection vector
\(w\).

\begin{table}[H]
\centering
\small
\begin{tabular}{ccccc}
\toprule
\(K\) & Reward & \(P(\mathrm{full})\) & Correct-action prob. & Missed@0.90 \\
\midrule
32 & 0.962 & 1.000 & 0.965 & 0.00 \\
\bottomrule
\end{tabular}
\caption{One-to-many selection-only experiment. Messages and policy are fixed,
and only the Observer's compressor-selection probabilities are learned.}
\label{tab:one-to-many-w-only}
\end{table}

\subsection{Policy Only: Learn \(\theta\), Fixed Correct Message and Full Selection}

This experiment fixes selection to the full compressor and fixes the message to
the correct scalar state value. The Observer only learns the policy
\(\pi_\theta(a\mid m)\).

\begin{table}[H]
\centering
\small
\begin{tabular}{cccccc}
\toprule
\(K\) & \(A\) & Reward & \(P(\mathrm{full})\) & Correct-action prob. & Missed@0.90 \\
\midrule
128 & 32 & 1.000 & 1.000 & 1.000 & 0.00 \\
\bottomrule
\end{tabular}
\caption{One-to-many policy-only experiment. Selection and messages are fixed,
and only the Observer policy is learned.}
\label{tab:one-to-many-theta-only}
\end{table}

\subsection{Selection and Compression: Learn \(\phi,w\), Fixed Policy}

This experiment keeps the Observer policy fixed to the distance-based softmax
policy from the one-to-one scalar-compression experiment. The learned objects are the selection logits
\(w\) and the selected compressor's message parameters \(\phi_{k_t}\). This
tests whether the Observer learns to listen to the full-information compressor
and whether that compressor learns the scalar sum message.

\reportfigure{updated_section_5_2_same_timescale_k16_eps020_alpha0003_300k_curves}
  {One-to-many learning at \(K=16\) with a fixed Observer policy. The same
  learning rate is used for selection and compression. The full-compressor
  probability moves close to one, and the scalar message becomes usable by the
  fixed policy.}
  {fig:section5-2-k16}

\reportfigure{updated_section_5_2_same_timescale_k16_eps020_alpha0003_300k_final_selection_messages}
  {One-to-many selection and compression at \(K=16\). The selection probability is
  concentrated on the full-information compressor, and the full compressor's
  scalar messages follow the state order.}
  {fig:section5-2-k16-final-messages}

\begin{table}[H]
\centering
\small
\resizebox{\linewidth}{!}{%
\begin{tabular}{cccccccc}
\toprule
\(K\) & Steps & lr & \(\epsilon_w\) & End expected & \(P(\mathrm{full})\) & Missed states & Success \\
\midrule
8  & 120k  & 0.010 & 0.05 & 0.967 & 0.996 & 0.00 & 1.00 \\
16 & 300k  & 0.003 & 0.20 & 0.963 & 0.997 & 0.00 & 1.00 \\
\bottomrule
\end{tabular}
}
\caption{One-to-many fixed-policy results over 5 seeds. Selection and
compression use the same learning rate.}
\label{tab:section5-2-scalar-sum}
\end{table}

Selection and compression work through \(K=16\) in this same-timescale setting.
The full-compressor probability moves close to one, and the selected full
compressor learns a scalar message that the fixed policy can use.

\subsection{Selection and Policy: Learn \(w,\theta\), Fixed Correct Messages}

This experiment fixes the compressor messages. The full compressor sends the
correct scalar state message, and the Observer learns both its selection
probabilities and its policy.

\begin{table}[H]
\centering
\small
\begin{tabular}{ccccccc}
\toprule
\(K\) & \(A\) & Reward & \(P(\mathrm{full})\) & Correct-action prob. & Msg MAE & Missed@0.90 \\
\midrule
128 & 32 & 0.971 & 0.973 & 1.000 & 0.000 & 0.00 \\
\bottomrule
\end{tabular}
\caption{One-to-many selection-and-policy experiment. Messages are fixed, while
selection and policy are learned.}
\label{tab:one-to-many-w-theta}
\end{table}

\subsection{Compression and Policy: Learn \(\phi,\theta\), Fixed Full Selection}

This experiment fixes selection to the full compressor and learns both the
scalar message and the policy from reward.

\begin{table}[H]
\centering
\small
\begin{tabular}{cccccc}
\toprule
\(K\) & \(A\) & Reward & Correct-action prob. & Msg MAE & Missed@0.90 \\
\midrule
128 & 32 & 0.999 & 0.983 & 0.145 & 0.00 \\
\bottomrule
\end{tabular}
\caption{One-to-many compression-and-policy experiment. Selection is fixed to
the full compressor.}
\label{tab:one-to-many-phi-theta}
\end{table}

\subsection{Full Learning: Learn \(\phi,w,\theta\)}

This experiment learns all three imperfect-state objects together: scalar
compression, compressor selection, and Observer policy.

\begin{table}[H]
\centering
\small
\begin{tabular}{ccccccc}
\toprule
\(K\) & \(A\) & Reward & \(P(\mathrm{full})\) & Correct-action prob. & Msg MAE & Missed@0.90 \\
\midrule
128 & 32 & 0.969 & 0.973 & 0.966 & 0.158 & 0.00 \\
\bottomrule
\end{tabular}
\caption{Full one-to-many imperfect-state experiment. Compression, selection,
and policy are learned together.}
\label{tab:one-to-many-full}
\end{table}

The one-to-many results show the intended ladder. Compression alone, selection
alone, and policy alone each learn. The two-object experiments also learn under
the fixed third object. The full experiment reaches high reward while preserving
the scalar state-message interpretation.

\section{Imperfect State with Multiple Actor-Compressors and Multiple Observers}

This section is the group imperfect-state setting in the updated paper. There
are three Observers and three Actor-compressors. Each Observer \(i\) has its
own selection logits \(w_i\) over compressors and its own policy
\(\pi_{\theta_i}(a\mid m)\). At each state, every Observer samples a compressor,
receives that compressor's scalar message, and acts. If several Observers
select the same compressor, that compressor receives feedback through all of
those selected paths.

\subsection{Compression Only: Learn \(\phi\), Fixed Full Selection and Fixed Policy}

This experiment fixes every Observer to select the full-information compressor
and fixes every Observer policy. The learned object is the full compressor's
message rule \(f_{\phi_2}\).

\begin{table}[H]
\centering
\small
\begin{tabular}{ccccc}
\toprule
\(K\) & Reward & \(P(\mathrm{full})\) & Correct-action prob. & Missed@0.90 \\
\midrule
32 & 0.963 & 1.000 & 0.966 & 0.00 \\
\bottomrule
\end{tabular}
\caption{Many-to-many compression-only experiment with fixed full selection and
fixed policies.}
\label{tab:many-to-many-phi-only}
\end{table}

\subsection{Selection Only: Learn \(w\), Fixed Correct Messages and Fixed Policy}

This experiment fixes the compressor messages and policies. Each Observer learns
only its compressor-selection probabilities.

\begin{table}[H]
\centering
\small
\begin{tabular}{ccccc}
\toprule
\(K\) & Reward & \(P(\mathrm{full})\) & Correct-action prob. & Missed@0.90 \\
\midrule
32 & 0.962 & 1.000 & 0.965 & 0.00 \\
\bottomrule
\end{tabular}
\caption{Many-to-many selection-only experiment. Messages and policies are
fixed, and each Observer learns which compressor to select.}
\label{tab:many-to-many-w-only}
\end{table}

\subsection{Policy Only: Learn \(\theta\), Fixed Correct Message and Full Selection}

This experiment fixes each Observer to receive the correct scalar message from
the full compressor. Each Observer only learns its policy.

\begin{table}[H]
\centering
\small
\begin{tabular}{cccccc}
\toprule
\(K\) & \(A\) & Reward & \(P(\mathrm{full})\) & Correct-action prob. & Missed@0.90 \\
\midrule
128 & 32 & 1.000 & 1.000 & 0.999 & 0.00 \\
\bottomrule
\end{tabular}
\caption{Many-to-many policy-only experiment. Messages and selection are fixed,
and the Observer policies are learned.}
\label{tab:many-to-many-theta-only}
\end{table}

\subsection{Selection and Compression: Learn \(\phi,w\), Fixed Policy}

This experiment fixes the Observer policies and learns both the scalar messages
and the per-Observer selection probabilities.

\begin{table}[H]
\centering
\small
\begin{tabular}{ccccc}
\toprule
\(K\) & Reward & \(P(\mathrm{full})\) & Correct-action prob. & Missed@0.90 \\
\midrule
32 & 0.961 & 1.000 & 0.964 & 0.00 \\
\bottomrule
\end{tabular}
\caption{Many-to-many selection-and-compression experiment with fixed policies.}
\label{tab:many-to-many-phi-w}
\end{table}

\subsection{Selection and Policy: Learn \(w,\theta\), Fixed Correct Messages}

This experiment fixes the compressor messages and learns each Observer's
selection probabilities and policy.

\begin{table}[H]
\centering
\small
\begin{tabular}{ccccccc}
\toprule
\(K\) & \(A\) & Reward & \(P(\mathrm{full})\) & Correct-action prob. & Msg MAE & Missed@0.90 \\
\midrule
16 & 16 & 0.995 & 0.999 & 0.994 & 0.000 & 0.00 \\
\bottomrule
\end{tabular}
\caption{Many-to-many selection-and-policy experiment. Messages are fixed, while
selection and policy are learned.}
\label{tab:many-to-many-w-theta}
\end{table}

\subsection{Compression and Policy: Learn \(\phi,\theta\), Fixed Full Selection}

This experiment fixes every Observer to use the full compressor and learns the
scalar messages and policies.

\begin{table}[H]
\centering
\small
\begin{tabular}{cccccc}
\toprule
\(K\) & \(A\) & Reward & Correct-action prob. & Msg MAE & Missed@0.90 \\
\midrule
128 & 32 & 0.999 & 0.986 & 0.130 & 0.00 \\
\bottomrule
\end{tabular}
\caption{Many-to-many compression-and-policy experiment with fixed full
selection.}
\label{tab:many-to-many-phi-theta}
\end{table}

\subsection{Full Learning: Learn \(\phi,w,\theta\)}

This experiment learns scalar compression, compressor selection, and Observer
policy together for all Observers. This run uses
\(\alpha_w=10^{-3}\) so that compressor selection settles within the 300k-step
training period; the other parameters use the common learned-policy
configuration given in the Method section.

\begin{table}[H]
\centering
\small
\begin{tabular}{ccccccc}
\toprule
\(K\) & \(A\) & Reward & \(P(\mathrm{full})\) & Correct-action prob. & Msg MAE & Missed@0.90 \\
\midrule
64 & 16 & 0.983 & 0.984 & 0.991 & 0.083 & 0.00 \\
\bottomrule
\end{tabular}
\caption{Full many-to-many imperfect-state experiment. Compression, selection,
and policy are learned together.}
\label{tab:many-to-many-full}
\end{table}

The many-to-many experiments show the same ladder as the one-to-many
experiments. Policy-only learning works from the correct scalar message,
compression plus policy works when selection is fixed to the full compressor,
and the full learned case works at \(K=64,A=16\).

\section{Discussion}

The experiments separate the two information problems in the paper. In the
imperfect-reward cases, the main learned object is the Observer's following
allocation and, in the joint cases, the Actor policy. In the imperfect-state
cases, the learned objects are the Actor-compressor message \(f_\phi\), the
Observer's compressor-selection probabilities \(\lambda\), and the Observer
policy \(\pi_\theta\).

The imperfect-reward experiments behave as expected in the controlled cases. With
fixed Actor quality, the following rates move toward the higher-utility Actors
and respect the budget. In the joint version, the Actors learn their policies
while the Observer's utility estimates and following rates move with them. The
group imperfect-reward case also shows the intended behavior: different
Observers can follow different Actors when their reward functions value
different Actors. The main scaling issue there is budget and sampling. When
many Actors are present, each good Actor must still receive enough active
updates and enough reported reward feedback.

The imperfect-state experiments use scalar-sum messages and a distance-based
softmax policy. The scalar-compression experiment works through \(K=64\), and
the one-Observer selection plus compression experiment works through \(K=16\)
in the same-timescale table. The full Actor-compressor learns a useful scalar
message, and the Observer learns to select it. The policy scores actions
according to distance from the scalar message to the action centers.

With this policy, the one-to-one and one-to-many imperfect-state cases work in
the reported \(K=128,A=32\) grouped-action setting. The many-to-many case works
at \(K=64,A=16\) when compression, selection, and policy are learned together.

The main results include reward and the internal learned quantities. For
imperfect-state learning, we need to record three learned objects separately:
whether the Actor-compressor sends useful state information, whether the
Observer can act from that message, and whether the compressor-selection
probabilities move toward the
full-information compressor. The tables include internal quantities such as message
error, message \(R^2\), correct-action probability, \(P(\mathrm{full})\), and
missed regions. The current conclusion is that the paper mechanism works in
the controlled subcases reported here, with the distance-based softmax policy
used for scalar messages.

\appendix
\section{Additional Experiments}

The appendix contains supporting experiments that are useful for interpreting the
main tables: following-rate timescales, tie cases, reward-baseline variants,
scaling experiments, and extra convergence plots.

Figure \ref{fig:appendix-paper4-timescale} compares three following-rate update
schedules on the same single-best controlled setup. The following-rate step is
held fixed at \(\alpha_\mu=0.02\), while the utility-estimate rate is varied.
When \(\alpha_U\) is larger, \(\hat U\) reaches the true utilities before the
following rates have moved far. When \(\alpha_U\) is equal to or smaller than
\(\alpha_\mu\), the same simple problem still converges, but the utility
estimate lags and the following-rate movement is correspondingly slower. This
supports the two-timescale convention used in the imperfect-reward
experiments: the utility estimate should track quickly, while \(\mu\) moves
more slowly under the budget projection.

\reportfigure{paper4_1_timescale_comparison}
  {Following-rate timescale comparison. The left panel shows mean utility-tracking
  error, and the right panel shows the following rate on the best Actor. Larger
  \(\alpha_U\) produces faster utility tracking before the slower
  following-rate update reallocates the attention budget.}
  {fig:appendix-paper4-timescale}

Figure \ref{fig:appendix-tie-cases} shows tie cases. In the action
tie, actions 0 and 1 both have reward \(1.0\), while action 2 has reward
\(0.0\). The exact split between actions 0 and 1 is seed-dependent, but the
low-reward action loses probability. In the Observer tie, two Observers provide
the same useful \(x_1+x_2\) message. The probability mass moves to the useful
set, even though the split between the two useful Observers can depend on
stochastic history.

\reportfigure{report_appendix_tie_cases}
  {Tie-case results. Tied optimal choices may split unevenly; the required
  behavior is that probability leaves the lower-value alternatives and
  concentrates on the optimal set.}
  {fig:appendix-tie-cases}

\subsection*{Raw Reward and Advantage Variant}

The paper's raw-reward selection update uses the raw reward multiplier
\[
  -r_t\log\lambda_{k_t}.
\]
We keep this as the main baseline. As an additional experiment, we compare it with
a global-advantage variant
\[
  -A_t\log\lambda_{k_t},
  \qquad
  A_t=r_t-b_t,
\]
where \(b_t\) is a global exponential moving average of reward. The baseline
depends only on the global reward history, so the change is the standard
variance-reduction form with the same credit assignment rule.

\reportfigure{paper5_1_advantage_positive_comparison}
  {Raw reward versus global advantage in the positive-reward
  setup. Both updates concentrate on the known best Observer. The advantage
  update changes the logit path and ends with almost the same end-of-run selection
  probability.}
  {fig:appendix-advantage-positive}

\reportfigure{paper5_1_advantage_summary}
  {Summary of the raw-reward and advantage comparison across
  selected conditions. The two updates are close in the original low-offset
  task, but the advantage update is more stable when rewards contain a large
  common positive offset.}
  {fig:appendix-advantage-summary}

\begin{table}[H]
\centering
\small
\resizebox{\linewidth}{!}{%
\begin{tabular}{llccccc}
\toprule
Condition & Update & Final reward & Final \(P(\text{best})\) &
Cumulative regret & Step to \(P(\text{best})>0.8\) & \(P(\text{best})\) std \\
\midrule
positive, gap \(1\) & raw reward & 0.958 & 0.995 & 116.1 & 231.0 & 0.001 \\
positive, gap \(1\) & advantage & 0.958 & 0.994 & 110.3 & 245.0 & 0.000 \\
offset \(5\), gap \(1\) & raw reward & 5.958 & 0.996 & 655.5 & 1149.4 & 0.002 \\
offset \(5\), gap \(1\) & advantage & 5.958 & 0.994 & 117.9 & 238.0 & 0.000 \\
offset \(20\), gap \(1\) & raw reward & 20.580 & 0.400 & 2097.5 & 208.8 & 0.490 \\
offset \(20\), gap \(1\) & advantage & 20.958 & 0.995 & 151.1 & 276.6 & 0.001 \\
offset \(5\), gap \(0.1\) & raw reward & 5.068 & 0.598 & 186.7 & 1394.2 & 0.486 \\
offset \(5\), gap \(0.1\) & advantage & 5.091 & 0.925 & 76.9 & 2298.0 & 0.011 \\
offset \(5\), gap \(1\), noise \(2\) & raw reward & 5.796 & 0.799 & 1005.7 & 1143.8 & 0.399 \\
offset \(5\), gap \(1\), noise \(2\) & advantage & 5.940 & 0.996 & 144.7 & 241.6 & 0.002 \\
quality hard & raw reward & 0.474 & 0.972 & 86.8 & 1125.4 & 0.005 \\
quality hard & advantage & 0.474 & 0.970 & 83.2 & 1105.0 & 0.006 \\
\bottomrule
\end{tabular}
}
\caption{Raw-reward versus global-advantage comparison over 5
seeds. The offset conditions add the same constant to all rewards, so the
best Observer is unchanged.}
\label{tab:appendix-advantage}
\end{table}

In the original low-offset selection task, the raw-reward and advantage
updates are nearly identical. The reward scale is small, the best Observer is
well separated, and the raw estimator has little variance to remove. The
difference appears when the reward table contains a large common positive
offset, such as rewards near \(20\) versus \(21\). That offset leaves the best
Observer unchanged and contributes little useful gradient in expectation, but it
multiplies the sampled log-probability and makes the raw update noisy.
The global advantage subtracts most of this common component. In the
offset-\(20\), gap-\(1\) test, raw reward ends with
\(P(\mathrm{best})=0.400\) and high seed variance, while the advantage update
ends with \(P(\mathrm{best})=0.995\). The step-to-threshold column can be
misleading in these offset runs because the raw update can briefly cross the
threshold and later move to a different Observer. The end-of-run best-Observer
probability, seed variance, and cumulative regret give the more useful
comparison.

The paper update should remain the raw-reward form because it is the simpler
baseline and matches the stated raw-reward rule. The advantage form is useful
as an appendix variance-reduction variant for reward designs with large common
offsets, small quality gaps, or high reward noise.

\subsection*{Imperfect-Reward Joint Scaling Experiments}
\label{appendix:imperfect-reward-joint-scaling}

This subsection folds in the imperfect-reward scaling experiments. The
setting is still the reward-following track. One Observer watches
many Actors. Two Actors are genuinely good: they see the state and can learn a
state-dependent policy. The other Actors are blind distractors: their input is
zeroed, so they can only learn one constant action. The Observer has following
budget \(B\), assigns following rates \(\mu_j\), and reports reward to Actor
\(j\) with probability
\[
  \theta(\mu_j)=1-e^{-c\mu_j}.
\]
A good Actor only takes a policy-gradient step when this report happens.

The earlier one-axis experiments were mostly stable. Scaling states alone worked.
Scaling the actor crowd alone mostly worked. The hard case was scaling states
and Actors together at a fixed budget. On a ladder from 16 states and 8 Actors
to 144 states and 128 Actors, good-Actor accuracy and budget concentration both
dropped sharply.

\reportfigure{scaling_states_actors_summary}
  {Joint state-and-Actor scaling at fixed budget \(B=2\). Good-Actor accuracy,
  budget concentration, and the concentration-weighted system score all drop
  sharply past the 64-state and 32-Actor rung.}
  {fig:appendix-joint-cliff}

\reportfigure{scaling_states_actors_curves}
  {Joint state-and-Actor scaling trajectories. The large-rung curves show the
  cliff throughout training and at the end of the run.}
  {fig:appendix-joint-cliff-curves}

\begin{table}[H]
\centering
\small
\begin{tabular}{ccccc}
\toprule
States & Actors & Good-Actor accuracy & Budget concentration & System score \\
\midrule
16  & 8   & 1.000 & 0.935 & 0.952 \\
36  & 16  & 0.917 & 0.839 & 0.816 \\
64  & 32  & 0.635 & 0.531 & 0.476 \\
100 & 64  & 0.302 & 0.052 & 0.258 \\
144 & 128 & 0.262 & 0.000 & 0.250 \\
\bottomrule
\end{tabular}
\caption{Joint imperfect-reward scaling at fixed \(B=2\), 60000 steps, and 3
seeds. Good-Actor accuracy is the actual policy-learning metric. Budget
concentration is the fraction of following budget on the two good Actors.}
\label{tab:appendix-ir-joint-cliff}
\end{table}

\textbf{Perfect-reward sanity experiment.} The first experiment hard-codes the report
probability to one. This means every acting Actor always receives its reward
feedback. It is the \(B\to\infty\) limit of the reporting rule, so it gives a
ceiling for what the imperfect-reward mechanism should approach when reporting
is frequent enough. With perfect reward, the joint cliff mostly disappears:
accuracy stays high where the fixed-budget run falls near chance.

\begin{table}[H]
\centering
\small
\begin{tabular}{rrcc}
\toprule
states & actors & perfect good-acc & \(B=2\) good-acc \\
\midrule
16  & 8   & 1.000 & 1.000 \\
36  & 16  & 0.986 & 0.907 \\
64  & 32  & 0.785 & 0.638 \\
100 & 64  & 0.792 & 0.302 \\
144 & 128 & 0.707 & 0.262 \\
\bottomrule
\end{tabular}
\caption{Perfect reward versus \(B=2\) along the joint ladder. Perfect reward
removes the sharp cliff, showing that the update rule has enough signal when
reward reports are frequent.}
\label{tab:appendix-ir-perfect-reward}
\end{table}

\textbf{Budget experiment.} At the collapsed 100-state and 64-Actor rung, increasing
\(B\) raises the report probability on good Actors and raises good-Actor
accuracy toward the perfect-reward ceiling. This also explains the earlier
confusion about concentration: concentration can fall when \(B\) is larger
because the budget is spread over more Actors, but good-Actor accuracy still
improves.

\reportfigure{reward_budget_sanity_summary}
  {Budget experiment at 100 states and 64 Actors. As \(B\) grows, report probability
  on the good Actors and good-Actor accuracy rise toward the perfect-reward
  line.}
  {fig:appendix-ir-budget-experiment}

\reportfigure{reward_budget_sanity_trajectories}
  {Budget experiment trajectories at 100 states and 64 Actors. Higher budget gives
  the good Actors more reward reports during training, and their accuracy
  improves with that extra feedback.}
  {fig:appendix-ir-budget-trajectories}

\begin{table}[H]
\centering
\small
\begin{tabular}{lcccc}
\toprule
condition & report prob. good & good-acc & concentration & system score \\
\midrule
perfect \((\theta=1)\) & 1.000 & 0.803 & 0.673 & 0.627 \\
\(B=2\)   & 0.048 & 0.302 & 0.052 & 0.258 \\
\(B=8\)   & 0.420 & 0.635 & 0.152 & 0.327 \\
\(B=32\)  & 0.703 & 0.788 & 0.077 & 0.292 \\
\(B=128\) & 0.940 & 0.793 & 0.044 & 0.274 \\
\bottomrule
\end{tabular}
\caption{Budget experiment at 100 states and 64 Actors. ``Report prob. good'' is the
mean reward-report probability for the two good Actors.}
\label{tab:appendix-ir-budget-experiment}
\end{table}

\textbf{Learning-rate experiment.} Peter's main instruction was to look at the
learning-rate curve before calling a case broken. Bo therefore swept the Actor
learning rate under perfect reward at the same 100-state and 64-Actor rung. The
curve has the expected shape: too small crawls, the middle range works best,
and too large locks into a bad policy.

\reportfigure{learning_rate_experiment_trajectories}
  {Good-Actor accuracy over training for each Actor learning rate under perfect
  reward. The middle rates climb; the largest rate moves badly and settles near
  chance.}
  {fig:appendix-ir-lr-trajectory}

\begin{table}[H]
\centering
\small
\begin{tabular}{rcc}
\toprule
learning rate & good-acc & late wobble \\
\midrule
0.001 & 0.623 & 0.028 \\
0.003 & 0.813 & 0.017 \\
0.010 & 0.803 & 0.035 \\
0.030 & 0.748 & 0.046 \\
0.100 & 0.250 & 0.000 \\
\bottomrule
\end{tabular}
\caption{Actor learning-rate experiment under perfect reward at 100 states and 64
Actors. Late wobble is the standard deviation of accuracy over late snapshots.}
\label{tab:appendix-ir-lr}
\end{table}

\textbf{Actor-selection dilution.} Even under perfect reward, uniform Actor
sampling gives each good Actor only a small share of the active updates in a
fixed-length run. With 64 Actors, the two good Actors together receive only 2
out of 64 turns under uniform sampling. Giving the good Actors more turn
probability closes most of the residual gap.

\reportfigure{actor_selection_experiment}
  {Good-Actor accuracy versus the selection mass placed on the good Actors.
  The dotted line is the uniform 2 out of 64 baseline. More turns per good Actor
  recover near-perfect accuracy.}
  {fig:appendix-ir-actor-selection}

\begin{table}[H]
\centering
\small
\begin{tabular}{rrc}
\toprule
good-Actor selection share & turns per good Actor & good-acc \\
\midrule
0.031 \;(uniform 2 out of 64) & 938   & 0.813 \\
0.10                         & 3000  & 0.968 \\
0.25                         & 7500  & 0.977 \\
0.50                         & 15000 & 0.962 \\
\bottomrule
\end{tabular}
\caption{Perfect reward, 100 states and 64 Actors, learning rate \(0.003\). The
residual gap is mostly the number of turns each good Actor gets.}
\label{tab:appendix-ir-selection-dilution}
\end{table}

\textbf{Putting it together.} With learning rate \(0.003\), budget \(B=32\),
and enough total steps, the once-collapsed rung works again under the ordinary
imperfect-reporting rule.

\begin{table}[H]
\centering
\small
\begin{tabular}{rccc}
\toprule
steps & report prob. good & good-acc & system score \\
\midrule
60{,}000  & 0.629 & 0.732 & 0.280 \\
300{,}000 & 0.743 & 0.930 & 0.308 \\
\bottomrule
\end{tabular}
\caption{Imperfect reward, 100 states and 64 Actors, learning rate \(0.003\),
\(B=32\). With enough compute the previously collapsed rung reaches 0.93
accuracy.}
\label{tab:appendix-ir-recovered-rung}
\end{table}

The conclusion is narrower than the first scaling read. The cliff came from
three controllable inputs: learning rate, reporting budget, and the number of
turns each good Actor receives. With those set correctly, the imperfect-reward
side recovers.

\textbf{Systematic per-axis stress test with scalar state ID and embedding.}
As an additional experiment under the representation used for the imperfect-reward runs,
Bo scaled one axis at a time. The state input is an integer ID
\(0,\ldots,K-1\). The policy first maps that ID to a learnable embedding vector,
then applies a linear policy head. Uniform action exploration anneals from
\(\varepsilon=0.1\) to \(0\) over training. A row is marked
\texttt{pass} when the settled optimal-action probability \(p_{\mathrm{opt}}\)
is near its ceiling, \texttt{watch} when it is lower but still usable, and
\texttt{rerun} when it needs more tuning.

\textbf{States.} This experiment uses one Actor, perfect reward, and \(K\) states
with \(K\) actions. It isolates the policy head on scalar IDs. A plain embedding
reaches \(p_{\mathrm{opt}}=1.0\) through \(K=16\). At \(K=32\) it drops to
\(0.833\), which is a \texttt{watch} row. The same embedding fits the diagnostic
lookup through much larger \(K\). This points to sampled policy-gradient
variance as the limiting factor. An MLP on the raw scalar learns less reliably
as \(K\) grows.

\begin{table}[H]
\centering
\small
\begin{tabular}{lcccc}
\toprule
policy head & \(K=3\) & \(K=8\) & \(K=16\) & \(K=32\) \\
\midrule
scalar plus embed       & 1.000 & 1.000 & 1.000 & 0.833 \\
scalar plus embed plus MLP1  & 1.000 & 1.000 & 0.750 & 0.271 \\
scalar plus rbf         & 1.000 & 1.000 & 0.937 & 0.687 \\
scalar plus MLP1, normalized & 0.889 & 0.458 & 0.187 & 0.052 \\
scalar plus MLP2, normalized & 0.778 & 0.167 & 0.083 & 0.031 \\
\bottomrule
\end{tabular}
\caption{States axis under scalar IDs: mean \(p_{\mathrm{opt}}\) over 3 seeds.
The plain embedding is the strongest policy head in this experiment.}
\label{tab:appendix-ir-stress-states}
\end{table}

\textbf{Actors.} This experiment uses one Observer, \(N\) Actors, and fixed
\(K=3\). The learned objects are the following rates \(\mu\) and the Actor
policies, as in the one-Observer imperfect-reward joint case. Mean and minimum
\(p_{\mathrm{opt}}\) stay above \(0.996\) through \(N=32\).

\begin{table}[H]
\centering
\small
\begin{tabular}{rccc}
\toprule
actors \(N\) & mean \(p_{\mathrm{opt}}\) & min \(p_{\mathrm{opt}}\) & verdict \\
\midrule
4  & 1.000 & 1.000 & pass \\
8  & 1.000 & 0.999 & pass \\
16 & 0.999 & 0.999 & pass \\
32 & 0.998 & 0.996 & pass \\
\bottomrule
\end{tabular}
\caption{Actor-axis stress test: one Observer and \(N\) Actors, scalar-ID
embedding policy, \(K=3\), 40000 steps, 3 seeds.}
\label{tab:appendix-ir-stress-actors}
\end{table}

\textbf{Observers.} This experiment fixes 8 Actors and \(K=3\), then varies the
number of Observers \(M\). All runs keep \(p_{\mathrm{opt}}=1.0\) through
\(M=32\).

\begin{table}[H]
\centering
\small
\begin{tabular}{rccc}
\toprule
observers \(M\) & mean \(p_{\mathrm{opt}}\) & min \(p_{\mathrm{opt}}\) & verdict \\
\midrule
4  & 1.000 & 1.000 & pass \\
8  & 1.000 & 1.000 & pass \\
16 & 1.000 & 1.000 & pass \\
32 & 1.000 & 1.000 & pass \\
\bottomrule
\end{tabular}
\caption{Observer-axis stress test: 8 Actors and \(M\) Observers, scalar-ID
embedding policy, \(K=3\), 40000 steps, 3 seeds.}
\label{tab:appendix-ir-stress-observers}
\end{table}

Taken together, under scalar state IDs with an embedding policy head, the
imperfect-reward system scales well when only states, only Actors, or only
Observers change. The only \texttt{watch} row is the large-\(K\) states-axis
policy run, and the earlier joint rung recovers once learning rate, budget, and
per-Actor compute are set correctly.

\subsection*{Imperfect-State Scaling Quantities}

The imperfect-state scaling dimensions reported in the main tables are:
\begin{itemize}
  \item number of states and actions \(K\),
  \item number of Observers and Actor-compressors,
  \item whether message parameters \(\phi\), selection logits \(w\), and policy
  parameters \(\theta\) are fixed or learned,
  \item the number of times each state and selected compressor receives useful
  feedback.
\end{itemize}

A missed region means a state and action region where the correct-action
probability is below the stated threshold, usually \(0.90\) or \(0.95\). This is
stricter than mean reward. It catches cases where average reward is high while
one state region still has low confidence.

The distance-based softmax policy uses the scalar message and supplies gradients
to the compression function. The main imperfect-state tables report
configurations whose reward and internal learned quantities reached the stated
thresholds.

\end{document}
